\subsection{Sigma-Algebras and Measurable Structures}

% **Role:** modern integration and almost-everywhere analysis.
%
% Sigma-algebras, measures, outer measure, measurable sets, Borel/Lebesgue measure, null sets, completion, charges, regularity.

Measure theory first builds measurable sets, then measurable functions and integrals. Outer measure covers from outside; measurability splits test sets without losing measure.

\begin{perspective} \label{per:measure-integration-roadmap} \hlt{Why Abstract Measure Comes First}\\
Lebesgue integration is built so that, under hypotheses such as domination, $f_n\to f$ implies $\int f_n\,d\mu\to\int f\,d\mu$; see Theorem \ref{thm:dominated-convergence-lebesgue}. The abstract data separate into a measurable space $(X,\mathcal{M})$, a countably additive measure $\mu:\mathcal{M}\to[0,+\infty]$, and an integral defined by approximation from simple functions.\\
\hlt{Key Insight}: Countable additivity is the set-level structure behind the convergence theorems.
\end{perspective}

\begin{definition} \label{def:sigma-algebra-and-measurable-space} \hlt{$\sigma$-Algebras and Measurable Spaces}\\
Let $X$ be a set. A family $\mathcal{M}\subseteq\mathcal{P}(X)$ is a \hlt{$\sigma$-algebra} on $X$ if:
\begin{enumerate}[label=(\roman*)]
\setlength{\itemsep}{0pt}
\item $\emptyset,X\in\mathcal{M}$;
\item $E\in\mathcal{M}$ implies $X\setminus E\in\mathcal{M}$;
\item $E_n\in\mathcal{M}$ for $n\in\N$ implies $\bigcup\limits_{n=1}^{\infty}E_n\in\mathcal{M}$.
\end{enumerate}
The pair $(X,\mathcal{M})$ is a \hlt{measurable space}; the members of $\mathcal{M}$ are \hlt{measurable sets}.
\end{definition}

\begin{proposition} \label{prop:sigma-algebra-closure-and-generation} \hlt{Closure and Generation for $\sigma$-Algebras}\\
Let $\mathcal{M}\subseteq\mathcal{P}(X)$ be a $\sigma$-algebra.
\begin{enumerate}[label=(\roman*)]
\setlength{\itemsep}{0pt}
\item $(E_n)_{n\in\N}\subseteq\mathcal{M}\Rightarrow\bigcap\limits_{n=1}^{\infty}E_n\in\mathcal{M}$.
\item If $(\mathcal{M}_\alpha)_{\alpha\in A}$ is a nonempty family of $\sigma$-algebras on $X$, then $\bigcap\limits_{\alpha\in A}\mathcal{M}_\alpha$ is a $\sigma$-algebra on $X$.
\item For $\mathcal{A}\subseteq\mathcal{P}(X)$, let $\mathscr{S}_{\mathcal{A}}$ be the family of $\sigma$-algebras $\mathcal{N}$ on $X$ with $\mathcal{A}\subseteq\mathcal{N}$.\\
Then $\sigma(\mathcal{A})\defeq\bigcap_{\mathcal{N}\in\mathscr{S}_{\mathcal{A}}}\mathcal{N}$ is the smallest $\sigma$-algebra containing $\mathcal{A}$.
\end{enumerate}
\begin{proof}
\begin{enumerate}[label=(\roman*)]
\setlength{\itemsep}{0pt}
\item Since $X\setminus E_n\in\mathcal{M}$, Theorem \ref{thm:metric-de-morgan-laws} gives $\bigcap_nE_n=X\setminus\bigcup_n(X\setminus E_n)\in\mathcal{M}$.
\item Set $\mathcal{N}\defeq\bigcap_{\alpha\in A}\mathcal{M}_\alpha$. Then $\emptyset,X\in\mathcal{N}$, and $E\in\mathcal{N}$ implies $X\setminus E\in\mathcal{N}$.\\
Likewise $(E_n)\subseteq\mathcal{N}$ implies $\bigcup_nE_n\in\mathcal{N}$.
\item The indexing class in the displayed definition is nonempty since $\mathcal{P}(X)$ belongs to it. By (ii), its intersection is a $\sigma$-algebra; it contains $\mathcal{A}$ and is contained in every $\sigma$-algebra containing $\mathcal{A}$.
\end{enumerate}
\hlt{Structural Commentary}: Generation is intersection over all admissible $\sigma$-algebras.\\
\hlt{Key Insight}: Intersections preserve universal closure properties.
\end{proof}
\end{proposition}

\begin{example} \label{ex:basic-sigma-algebras-and-borel} \hlt{Basic $\sigma$-Algebras}\\
Standard models:
\begin{enumerate}[label=(\roman*)]
\setlength{\itemsep}{0pt}
\item Extremes: $\{\emptyset,X\}\subseteq\mathcal{M}\subseteq\mathcal{P}(X)$ for every $\sigma$-algebra $\mathcal{M}$ on $X$.
\item For $E\subseteq X$, $\sigma(\{E\})=\{\emptyset,E,X\setminus E,X\}$, with repetitions omitted.
\item On $\N$, if $O\defeq\{2n-1:n\in\N\}$, then $\sigma(\{O\})=\{\emptyset,O,\N\setminus O,\N\}$.
\item If $X$ is uncountable, $\mathcal{C}_X\defeq\{E\subseteq X:E\text{ is countable or }X\setminus E\text{ is countable}\}$\\
is the countable--cocountable $\sigma$-algebra.
\end{enumerate}
\hlt{Key Insight}: A $\sigma$-algebra records exactly which distinctions the theory is allowed to observe.
\end{example}

\begin{definition} \label{def:borel-and-extended-borel-sigma-algebras} \hlt{Borel and Extended Borel $\sigma$-Algebras}
\begin{enumerate}[label=(\roman*)]
\setlength{\itemsep}{0pt}
\item For $n\in\N$, the \hlt{Borel $\sigma$-algebra} on $\R^n$ is $\mathcal{B}(\R^n)\defeq\sigma(\{U\subseteq\R^n:U\text{ is open}\})$.\\
On $\R$, it is also generated by open intervals, closed intervals, half-open intervals, or open half-rays $(a,+\infty)$.
\item The \hlt{extended Borel $\sigma$-algebra} on $\overline{\R}$ is $\overline{\mathcal{B}}\defeq\{E\cup F:E\in\mathcal{B}(\R),\ F\subseteq\{-\infty,+\infty\}\}$.
\end{enumerate}
\end{definition}

\begin{definition} \label{def:monotone-classes} \hlt{Monotone Classes}\\
A nonempty $\mathcal{C}\subseteq\mathcal{P}(X)$ is a \hlt{monotone class} iff $(E_n)\subseteq\mathcal{C},\ E_n\uparrow E\Rightarrow E\in\mathcal{C}$ and $(F_n)\subseteq\mathcal{C},\ F_n\downarrow F\Rightarrow F\in\mathcal{C}$.\\
For $\mathcal{A}\subseteq\mathcal{P}(X)$, define $\mathfrak{M}_{\mathcal{A}}\defeq \{\mathcal{C}\subseteq\mathcal{P}(X):\mathcal{C}\text{ is a monotone class},\ \mathcal{A}\subseteq\mathcal{C}\}$; $\operatorname{mon}(\mathcal{A})\defeq\bigcap\mathfrak{M}_{\mathcal{A}}$.\\
Then $\operatorname{mon}(\mathcal{A})$ is the least monotone class containing $\mathcal{A}$, and $\operatorname{mon}(\mathcal{A})\subseteq\sigma(\mathcal{A})$.\\
\hlt{Key Insight}: Monotone closure captures limits, not finite unions.
\end{definition}

\begin{definition} \label{def:algebra-field-and-measure-on-algebra} \hlt{Algebras of Sets and Premeasures}\\
Let $X$ be a set.
\begin{enumerate}[label=(\roman*)]
\setlength{\itemsep}{0pt}
\item An \hlt{algebra of sets} (or \hlt{field of sets}) on $X$ is $\mathcal{A}\subseteq\mathcal{P}(X)$ such that $\emptyset,X\in\mathcal{A}$; $E\in\mathcal{A}\Rightarrow X\setminus E\in\mathcal{A}$;\\
and $E_1,\ldots,E_n\in\mathcal{A}\Rightarrow\bigcup_{j=1}^{n}E_j\in\mathcal{A}$. Equivalently, $\mathcal{A}$ is a ring of sets containing $X$.
\item A \hlt{premeasure} (or \hlt{measure on an algebra}) on $\mathcal{A}$ is $\mu:\mathcal{A}\to[0,+\infty]$ such that $\mu(\emptyset)=0$ and, whenever $(E_n)_{n\in\N}\subseteq\mathcal{A}$ is pairwise disjoint with $\bigcup_{n=1}^{\infty}E_n\in\mathcal{A}$, then $\mu\left(\bigcup_{n=1}^{\infty}E_n\right)=\sum_{n=1}^{\infty}\mu(E_n)$.
\end{enumerate}
\hlt{Key Insight}: Premeasures are countably additive where the algebra permits unions.
\end{definition}

\begin{theorem} \label{thm:monotone-class-theorem-algebras} \hlt{Monotone Class Theorem for Algebras}\\
Let $\mathcal{A}\subseteq\mathcal{P}(X)$ be an algebra and set $\mathcal{C}\defeq\operatorname{mon}(\mathcal{A})$. Then $\mathcal{C}=\sigma(\mathcal{A})$.
\begin{proof} Since $\sigma(\mathcal{A})$ is monotone and contains $\mathcal{A}$, $\mathcal{C}\subseteq\sigma(\mathcal{A})$. It suffices to show that $\mathcal{C}$ is a $\sigma$-algebra.\\
For complements, set $\mathcal{D}\defeq\{E\in\mathcal{C}:E^c\in\mathcal{C}\}$. Then $E_n\uparrow E$ in $\mathcal{D}$ implies $E_n^c\downarrow E^c$, and $E_n\downarrow E$ implies $E_n^c\uparrow E^c$; hence $\mathcal{D}$ is monotone. Since $\mathcal{A}\subseteq\mathcal{D}$, minimality gives $\mathcal{C}\subseteq\mathcal{D}$.\\
For intersections, fix $B\in\mathcal{C}$ and set $\mathcal{D}_B\defeq\{E\in\mathcal{C}:E\cap B\in\mathcal{C}\}$. This is monotone because intersections preserve $\uparrow,\downarrow$ limits. If $B\in\mathcal{A}$, then $\mathcal{A}\subseteq\mathcal{D}_B$, so $\mathcal{C}\subseteq\mathcal{D}_B$; hence $E\cap B\in\mathcal{C}$ for $E\in\mathcal{C}$, $B\in\mathcal{A}$. Fixing $E\in\mathcal{C}$ now gives $\mathcal{A}\subseteq\mathcal{D}_E$, hence $\mathcal{C}\subseteq\mathcal{D}_E$. Thus $\mathcal{C}$ is closed under finite intersections.\\
Complements and finite intersections give finite unions. If $(E_n)\subseteq\mathcal{C}$, then $F_N\defeq\bigcup_{n=1}^{N}E_n\in\mathcal{C}$ and $F_N\uparrow\bigcup_nE_n$, so $\bigcup_nE_n\in\mathcal{C}$. Hence $\mathcal{C}$ is a $\sigma$-algebra containing $\mathcal{A}$, and $\sigma(\mathcal{A})\subseteq\mathcal{C}$.\\
\hlt{Structural Commentary}: Finite Boolean closure plus monotone limits gives $\sigma$-closure.\\
\hlt{Key Insight}: Prove closure properties by showing the sets that have the property form a monotone class.
\end{proof}
\end{theorem}

\begin{corollary} \label{cor:monotone-class-algebra-test} \hlt{Monotone-Class Test}\\
Let $\mathcal{A}\subseteq\mathcal{P}(X)$ be an algebra and let $\mathcal{C}\subseteq\mathcal{P}(X)$ be a monotone class. If $\mathcal{A}\subseteq\mathcal{C}$, then $\sigma(\mathcal{A})\subseteq\mathcal{C}$.
\begin{proof} Since $\mathcal{A}\subseteq\mathcal{C}$ and $\mathcal{C}$ is monotone, $\operatorname{mon}(\mathcal{A})\subseteq\mathcal{C}$. By Theorem \ref{thm:monotone-class-theorem-algebras}, $\operatorname{mon}(\mathcal{A})=\sigma(\mathcal{A})$.\\
\hlt{Structural Commentary}: Algebra checks plus monotone stability control $\sigma(\mathcal{A})$.\\
\hlt{Key Insight}: Prove on the algebra; extend by monotone closure.
\end{proof}
\end{corollary}
